\subsection{The degree-five discriminant and boundary periods}\label{sec03:degree5}

The degree-five cone has a smooth local deformation base, so it will
remain unchanged on the global partial resolution.  Its discriminant
requires five separate nonvanishing conditions.

\begin{proposition}\label{prop03:degree5family}
The full semiuniversal deformation base of $C_5$ is a smooth
four-dimensional germ.  It admits linear coordinates in which the
reduced discriminant is the union of five distinct hyperplanes.
Every noncentral singular fibre has only ordinary double points,
one for each discriminant hyperplane containing its parameter.
\end{proposition}

\begin{proof}
Let $V=\CC^5$, $A=\bigwedge^2V$, and
\[
 D=\{p\in A:p\wedge p=0\}=\Cone(\operatorname{Gr}(2,V)).
\]
Choose a general six-dimensional subspace $W\subset A$.
The surface $E=\operatorname{Gr}(2,V)\cap\PP(W)$ is smooth by
Bertini.  Adjunction gives $K_E=-H$, and its degree is five.
All smooth degree-five del Pezzo surfaces over $\CC$ are isomorphic,
so its anticanonical cone is the required germ.  Put
$U=W^\perp\subset A^*$ and $B=A/W=U^*$.  The quotient map
\begin{equation}\label{eq03:degree5family}
 L:D\longrightarrow B
\end{equation}
varies the four cutting linear forms by constants and has central
fibre $C_5$.

The Grassmannian cone is Cohen--Macaulay of dimension seven.
For example, its five quadratic Pfaffians have the graded resolution
\[
 0\longrightarrow S(-5)\longrightarrow S(-3)^5
 \longrightarrow S(-2)^5\longrightarrow S
 \longrightarrow\cO_D\longrightarrow0,
 \qquad S=\operatorname{Sym}(A^*).
\]
The four components of $L$ cut a three-dimensional central fibre
and hence form a regular sequence at the vertex.  Miracle flatness
over the smooth four-dimensional base proves flatness there.
Homogeneity then proves flatness everywhere: scaling transports
each point into this neighborhood and acts compatibly on $B$.

We check the full Kodaira--Spencer map.  A constant vector $c\in A$
gives the perturbation $p\wedge c$ of the Pfaffian equations on $W$,
up to the common factor two.  If this abstract first-order class is
zero, the cotangent presentation of the embedded germ shows that
its homogeneous degree-$-1$ part comes from a constant translation
$w\in W$.  The quadratic ideal has no linear elements, so
\[
 p\wedge(c-w)=0\quad\hbox{for every }p\in W.
\]
Write $a=c-w$.  If $a$ has rank four, the map
$a\wedge-:\bigwedge^2V\to\bigwedge^4V$ has rank five; its
kernel cannot contain the six-dimensional space $W$.
If $a$ has rank two with support $P$, its kernel is
$P\wedge V$, of dimension seven.  The decomposable bivectors in
$\PP(P\wedge V)$ form the four-dimensional Schubert variety of
two-planes meeting $P$: choose a line in $P$ and then a second
direction modulo that line.  A hyperplane $\PP(W)$ in this
$\PP^6$ would meet the Schubert variety in dimension at least
three, contradicting $\dim E=2$.  Thus $a=0$.
The Kodaira--Spencer map $B\to T^1_{C_5}$ is injective and is
an isomorphism because $\dim T^1_{C_5}=4$
\cite[\S5.5]{Gross97rev}.  This also shows that all of $T^1_{C_5}$
has degree $-1$.

Embed a semiuniversal base minimally with ring
$\CC\{z_1,\ldots,z_4\}/I$, $I\subset(z_1,\ldots,z_4)^2$.
The classifying map of~\eqref{eq03:degree5family} pulls the $z_i$
back to four functions with independent linear terms.  They are
analytic coordinates on $B$.  Hence the map from the ambient
convergent power-series ring is injective, forcing $I=0$.
The family is therefore semiuniversal for the full germ, including
its scheme structure; local prorepresentability is not required.

For a nonzero decomposable $p$ with support $P_p$, one has
$T_pD=P_p\wedge V$.  An alternating form $\alpha\in A^*$
annihilates this tangent space precisely when
$P_p\subset\ker\alpha$.  Such a nonzero form has rank two.
Thus the dual Grassmannian is the rank-two locus in $\PP(A^*)$.
It meets $\PP(U)$ transversely in five points
$[\alpha_1],\ldots,[\alpha_5]$.  Indeed, for a rank-two
$\alpha\in U$, failure of transversality is equivalent to
\[
 W\cap\bigwedge^2\ker\alpha\ne0.
\]
Every nonzero bivector in this intersection is decomposable and
would be a nonvertex singular point of $C_5$, contrary to the
smoothness of $E$.  The same argument excludes a positive-dimensional
intersection.  Its length is the degree five of the Grassmannian.

Set $K_i=\ker\alpha_i$ and $H_i=\ker(\alpha_i:B\to\CC)$.
The critical locus of $L$ away from the vertex is the union of
the punctured three-dimensional vector spaces $\bigwedge^2K_i$.
The maps
\[
 L:\bigwedge^2K_i\longrightarrow H_i
\]
are isomorphisms: their kernels are the intersections just excluded,
and both dimensions are three.  Thus a point $b\in H_i\setminus0$
has one associated critical point $p_i(b)$, linear in $b$.
Two such points cannot coincide for $b\ne0$.  Otherwise every
form in the pencil spanned by $\alpha_i,\alpha_j$ would vanish
on the same two-plane and hence have rank at most two, giving a
line in the finite dual intersection.  The complete reduced
discriminant is consequently $\bigcup_iH_i$.

Finally, near a nonzero $p_i(b)$ choose a Grassmannian chart in
which $\alpha_i=e_4^*\wedge e_5^*$ and the varying plane is
spanned by
\[
 e_1+x_3e_3+x_4e_4+x_5e_5,\qquad
 e_2+y_3e_3+y_4e_4+y_5e_5.
\]
With the nonzero cone coordinate $r$, the function $\alpha_i$
is $r(x_4y_5-x_5y_4)$.  Three other components of $L$ restrict
to coordinates on its critical locus, by the preceding isomorphism.
Eliminate them.  The remaining equation has nondegenerate quadratic
part in $x_4,x_5,y_4,y_5$, so the holomorphic Morse lemma gives
an ordinary double point with transverse smoothing parameter
$\alpha_i(b)$, up to a unit.
\end{proof}

\begin{lemma}\label{lem03:degree5marking}
Under the local Hodge comparison and the conic-pencil coordinates
\eqref{eq:degree5coordinates}, the discriminant of
Proposition~\ref{prop03:degree5family} is
\begin{equation}\label{eq03:degree5discriminant}
 \prod_{i=1}^5\lambda_i=0,
 \qquad \sum_{i=1}^5\lambda_i=0.
\end{equation}
\end{lemma}

\begin{proof}
Permuting the four blown-up points gives the $S_4$ action on the
first four conic pencils.  The quadratic Cremona involution at
the first three points lifts to an automorphism of $E$ and
interchanges the fourth and fifth pencils, fixing the first three.
These automorphisms generate $S_5$.  On
$\operatorname{Tot}(K_E)$ their canonical lifts preserve the
tautological two-form $\eta$ and the volume form $d\eta$.
The local Hodge comparison is therefore equivariant and realizes
$T^1_{C_5}$ as the standard sum-zero representation on the five
coordinates $\lambda_i$.

The tangent cone of the discriminant is intrinsic to the
deformation functor.  An automorphism of the central germ induces
a classifying map on a semiuniversal base; its invertible derivative
is the natural action on $T^1$.  It preserves the discriminant
because it preserves the isomorphism classes of the fibres.
Hence $S_5$ permutes the five hyperplanes in
Proposition~\ref{prop03:degree5family}.  This permutation action
is faithful.  A nontrivial kernel would contain $A_5$ and make
each normal line invariant under $A_5$.  As $A_5$ has no
nontrivial characters and no invariants in the standard
four-dimensional representation, this is impossible.

The image of this faithful action has order $120$, so it is the
full symmetric group on the five hyperplanes.  The stabilizer
of one hyperplane has order $24$.  In its action
on the five conic pencils it must fix a pencil: without a fixed
point its orbits would have sizes $2+3$, giving order at most
$2!3!=12$, while a size-five orbit cannot divide $24$.
It is thus the $S_4$ stabilizing one pencil.  The restriction
of the standard representation to this subgroup is the sum
of a trivial line and its three-dimensional standard representation.
The unique invariant normal line is generated by the corresponding
$\lambda_i$.  This identifies all five hyperplanes.

For a global volume form $g\,d\eta$, multiplication by the
analytic unit $g$ acts on $T^1_{C_5}$ by $g(0)$: its degree-$-1$
concentration gives $\mathfrak m_0T^1_{C_5}=0$.
Thus the fixed global normalization changes the comparison by
one nonzero scalar and preserves the five hyperplanes.
\end{proof}

\begin{lemma}\label{lem03:degree5milnor}
For a smooth Milnor fibre $M$ of $C_5$ with boundary link $L$,
\[
 H_2(M;\QQ)=0,\qquad \dim H_3(M;\QQ)=4.
\]
Both maps
\begin{equation}\label{eq03:degree5boundary}
 H_3(L;\QQ)\longrightarrow H_3(M;\QQ),\qquad
 H_3(M,L;\QQ)\longrightarrow H_2(L;\QQ)
\end{equation}
are isomorphisms.
\end{lemma}

\begin{proof}
For $b$ outside the discriminant the affine fibre compactifies as
\[
 V_b=\operatorname{Gr}(2,V)\cap\PP(L^{-1}(\CC b)),
 \qquad X_b=V_b\setminus E.
\]
This is a smooth three-hyperplane section of the Grassmannian.
It is smooth at $E$ because the four-hyperplane section $E$ is
smooth, and away from $E$ by Proposition~\ref{prop03:degree5family}.
Near the compact smooth divisor $E$, the squared affine norm is
$F/|s|^2$, where $s$ is a local equation for $E$ and $F$ is
smooth and positive.  Its normal derivative does not vanish
sufficiently close to $E$.  Thus every sufficiently large
Euclidean truncation of $X_b$ removes only a collar at infinity
and has its homotopy type.  Scaling by $\varepsilon$ identifies
$X_b\cap\{\|p\|\le R\}$ with
$X_{\varepsilon b}\cap\{\|p\|\le\varepsilon R\}$.
Fix the latter radius and take $R$ large.  This gives a small
Milnor fibre; transport in the smooth local family gives the
assertion for every nearby smooth parameter.

Lefschetz gives $b_1(V_b)=0$ and $b_2(V_b)=1$.
For completeness the middle Betti number follows from the Chern
calculation.  Write $H=c_1(\mathcal S^*)$ and
$P=c_2(\mathcal S^*)$ for the tautological bundle on the
Grassmannian.  Its tangent bundle is $\mathcal S^*\otimes\mathcal Q$,
so
\[
 c_1=5H,\qquad c_2=11H^2+P,\qquad c_3=15H^3.
\]
Dividing by $(1+H)^3$ gives $c_3(TV_b)=2H^3-3HP$.
Pieri's rule gives
$\int_{V_b}H^3=5$ and $\int_{V_b}HP=2$; hence
$\chi(V_b)=4$ and $b_3(V_b)=0$.
The Gysin sequence for $E\subset V_b$ gives
$H^2(X_b;\QQ)=0$ and
\[
 H^3(X_b;\QQ)\cong
 \ker\{H^2(E;\QQ)\longrightarrow H^4(V_b;\QQ)\},
\]
of dimension four.  The link circle-bundle sequence likewise
gives $b_2(L)=b_3(L)=4$.  Lefschetz duality and the pair sequence
make the second map in~\eqref{eq03:degree5boundary} a surjection
between four-dimensional spaces.  It is an isomorphism.  Exactness
then makes the first map surjective and hence an isomorphism as well.
\end{proof}

\begin{lemma}[degree-five boundary periods]\label{lem03:degree5periods}
There are rational classes $\beta_1,\ldots,\beta_5\in H_3(L;\QQ)$
whose pairings with the local Hodge classes are nonzero multiples
of the five functionals $\lambda_i$.
For any convergent smoothing arc of $C_5$, equipped with a relative
volume form generating its dualizing sheaf, the periods of the
images of all five classes in the smooth fibre are nonzero for
every sufficiently small noncentral parameter.
\end{lemma}

\begin{proof}
Use the semiuniversal family~\eqref{eq03:degree5family}.
The boundary family of a small representative is a smooth
proper submersion over the entire small base $B$, including $0$;
its homology local system is therefore trivial.
By~\eqref{eq03:degree5boundary}, every absolute Milnor cycle
has a unique boundary class.  Near a general point of the
discriminant hyperplane $H_i$, choose the class $\beta_i$
whose image is the nodal vanishing sphere, and transport it in
the boundary family over $B$.

The Pfaffian resolution above shows that the Grassmannian cone
has a homogeneous dualizing generator of weight five.  Dividing
by the four base differentials gives a nowhere vanishing
relative generator $\Omega_b^0$ of weight one.
Integration on transported compact boundary cycles gives
holomorphic functions on the whole base,
\[
 P_i(b)=\int_{\beta_{i,b}}\Omega_b^0.
\]
Holomorphic dependence uses only the smooth locus near the
boundary, so it holds also at parameters with a singular interior.
Scaling changes the boundary radius; compare the two boundaries
through their annular collars, whose central radial isotopy fixes
the homology marking.  The connected scaling action then preserves
transported classes and gives $P_i(cb)=cP_i(b)$.  Thus $P_i$ is linear.
Near a general point of $H_i$, it equals the nodal vanishing
period and has a simple zero along $H_i$.  Consequently
\begin{equation}\label{eq03:degree5linearperiod}
 P_i(b)=c_i\alpha_i(b),\qquad c_i\ne0.
\end{equation}

We verify that its derivative is the claimed Hodge functional.
Choose first-order product trivializations on the punctured
central germ, with Kodaira--Spencer cocycle $\theta_{jk}$.
Differentiating the relative three-form gives
\[
 \dot\Omega_k-\dot\Omega_j
       =d(\iota_{\theta_{jk}}\Omega_0^0).
\]
The \v{C}ech--de Rham representative for the derivative of a
period therefore has overlap term
$\iota_{\theta_{jk}}\Omega_0^0$, precisely the contraction
cocycle defining Lemma~\ref{lem:localhodge}.
In the logarithmic de Rham model of the circle bundle, this is
the natural comparison underlying \cite[Theorem~2.1(vi)]{FL22a}.
The link has $H^3(L;\CC)$ of pure Hodge type $(2,2)$, so its
$F^3$ is zero and this comparison determines the entire class.
Thus $dP_i|_0$ pairs $\beta_i$ with the local Hodge class,
up to the one common convention for Gysin maps and orientation.
There is no ambiguity from the chosen extension of the volume
form: a difference with the same overlap cocycle is a global
form $a\Omega_0^0$.  Normality extends $a$ across the vertex,
and contraction with the Euler vector field makes this form
exact, since all its homogeneous weights are strictly positive.
The same identity applies to a convergent series after shrinking.
Lemma~\ref{lem03:degree5marking} and
\eqref{eq03:degree5linearperiod} identify its kernel with
$\lambda_i=0$.

Finally let $b=b(t)$ be an arbitrary smoothing arc.  Its actual
relative volume form is $g(p,t)\Omega^0$, where $g$ is an
analytic unit on the pulled-back family.  Lift $g$ to a holomorphic
function of the ambient coordinates $(p,t)$; after shrinking it
is a unit and defines the same multiplier over $B\times\Delta$.
The boundary periods are holomorphic functions $P_i^g(b,t)$.
They vanish on $H_i\times\Delta$, by the nodal vanishing-cycle
calculation at its general points and analytic continuation.
Hence
\[
 P_i^g(b,t)=\alpha_i(b)u_i(b,t).
\]
To compute its linear term without expanding $g$ term by term,
fix a smooth direction $b_*$ and a compact cycle $\delta$
representing $\beta_i$ in $X_{b_*}$.  Scaling this cycle gives
\[
 P_i^g(sb_*,t)
 =s\int_\delta g(sp,t)\Omega_{b_*}^0
 =s\,g(0,t)c_i\alpha_i(b_*)+O(s^2).
\]
The estimate is uniform on the fixed compact cycle.  Smooth
directions form a dense open set, so
$d_bP_i^g(0,t)=g(0,t)c_i\alpha_i$.  Thus
$u_i(0,0)=g(0,0)c_i\ne0$.  Restricting to $b(t)$ proves
nonvanishing, whatever its contact order with $H_i$.
\end{proof}
